\begin{tikzpicture}[font=\small,>=stealth,line join=round]
\draw[draw=teal!45!black,fill=green!3,rounded corners=3pt] (0,0) rectangle (13.6,5.9);
\node at (6.8,5.3500000000000005) {反推的关键：$H$、$K$ 的导数为零，各自图像都是水平线};
\begin{scope}[shift={(4.8,1.3)},x=1.35cm,y=0.18cm]
\draw[gray,->] (-2.6,0)--(2.7,0) node[right] {$t$};
\draw[gray,->] (0,-1)--(0,17);
\draw[teal!75!black,very thick] (-2.3,0)--(2.3,0);
\draw[blue!70!black,very thick] (-2.3,14)--(2.3,14);
\draw[gray,dashed] (0,0)--(0,14);
\fill[teal!75!black] (0,0) circle(2.5pt);
\fill[blue!70!black] (0,14) circle(2.5pt);
\node[below] at (0,-.2) {$0$};
\end{scope}
\node[anchor=west,align=left,text width=3.8cm] at (9,3.8) {$K(t)\equiv14$\\取 $f(x)=(x-1)^3+7$\\$a=1,\ K(0)=2f(1)=14$。};
\node[anchor=west,align=left,text width=3.8cm] at (9,1.6) {$H(t)\equiv0$\\取 $f(x)=(x-2)^2+1$\\$a=2,\ H(0)=0$。};
\node at (6.8,.4) {图示用具体函数展示；正文证明对任意符合条件的 $f$ 成立。};
\end{tikzpicture}
